\section{Related Work}
\label{sec:related}

\textbf{State representation.}
Modern LIO includes iterated filters, smoothers, and continuous-time
estimators. FAST-LIO/FAST-LIO2 place gravity on $S^2$ inside an iterated
error-state filter \cite{xu2021fastlio,xu2022fastlio2,he2021ikfom}, and
Point-LIO retains the same convention \cite{he2023pointlio}. VE-LIOM also
estimates gravity online in an optimization framework \cite{gao2024veliom},
whereas LIO-SAM initializes a gravity-aligned frame without an equivalent
online direction state \cite{shan2020liosam}. Related pipelines also differ in
feature construction and map access. LOAM introduced edge/plane registration
\cite{zhang2014loam}, while FAST-LIO2 couples raw points to an incremental
$k$-d tree \cite{cai2021ikdtree}. Direct and continuous-time systems change
scan matching and propagation together
\cite{wang2023dliom,nguyen2023slict,chen2023dlio}. Cross-system accuracy
therefore confounds gravity handling with the front end, trajectory model,
map, and parameter settings. Our primary intervention changes only the compiled state
manifold; LIO-SAM provides a descriptive transfer check.

\textbf{Gravity--attitude--bias observability.}
Gravity direction, roll/pitch, and accelerometer bias have coupled
linearizations. Visual--inertial consistency analyses identify unobservable
and weakly observable directions \cite{hesch2014consistency,qin2018vinsmono},
and planar motion further limits excitation \cite{wu2017vinswheels}. LIO
initializers consequently estimate gravity, biases, timing, and extrinsics
before normal operation \cite{zhu2022liinit}; joint on-manifold calibration
makes the coupling explicit \cite{nemiroff2023gravity}. These analyses
establish when gravity is identifiable, but not whether its continued freedom
improves pose under frequent LiDAR correction. We test the latter while
intervening independently on $\mathbf g$ and $\mathbf b_a$. Mature
visual--inertial systems combine observability decisions with initialization,
relocalization, and map management \cite{campos2021orbslam3}; their aggregate
robustness cannot isolate one state's post-initialization utility.

\textbf{Gravity-enhanced residuals.}
Gravity-constrained registration removes rotational freedom using an
IMU-derived vertical direction \cite{kubelka2022gravity}. Recent radar--LiDAR
and radar--leg estimators add velocity-supported or soft $S^2$ gravity
information and report improved vertical accuracy
\cite{noh2025garlio,noh2025garlileo}. Elevator-specific models likewise show
that non-inertial motion can violate a nominal gravity model
\cite{elevator2026lio}. These methods add sensor information or residual
constraints; fixing an initialized gravity state does neither. We therefore
cross a same-IMU gravity-direction factor with fixed/online gravity and pair
its residual reduction with trajectory error. This design tests whether
tighter down-direction consistency reliably reduces vertical drift.

\textbf{Degradation and correction absence.}
Perceptual degeneracy is commonly detected from LiDAR geometry or estimator
observability and handled through selective updates, alternate odometry, or
direction-dependent weighting \cite{tagliabue2021lion,yao2025d2lio}. Weak
geometry and missing correction are often discussed within the same robustness
setting, although their estimator inputs differ. We separate them
experimentally: range/FoV degradation preserves the scan-update rhythm,
whereas periodic dropout removes correction entirely. Because the expected
effects are small, our protocol also audits reference-path plausibility and
exact correction timestamps, consistent with uncertainty-aware benchmark
generation \cite{hu2024paloc}. Our claim concerns an observation regime, not a
ranking of complete estimators. It distinguishes weak but accepted updates
from rejected or unavailable ones, a difference hidden by a single trajectory
score.
